\documentclass[a4paper,12pt]{article}
\usepackage[utf8]{inputenc}
\usepackage[T1]{fontenc}
\usepackage[french]{babel}
\usepackage{geometry}
\usepackage{xcolor}
\usepackage[explicit]{titlesec}
\usepackage{tcolorbox}
\usepackage{lmodern}
\usepackage{enumitem}
\usepackage{amssymb} % Nécessaire pour \mathbb, \forall, \exists

\geometry{margin=2.5cm}

% --- Style des sections ---
\titleformat{\section}
{\normalfont\Large\bfseries\sffamily}
{}
{0pt}
{\begin{tcolorbox}[colback=blue!5!white, colframe=blue!75!black, left=5pt, sharp corners, boxrule=0pt, leftrule=3pt] \thesection. #1 \end{tcolorbox}}

% --- Style des listes (itemize/enumerate) ---
\setlist[itemize,1]{label=\color{blue!75!black}\bullet, font=\bfseries, leftmargin=2em, labelsep=1em}
\setlist[enumerate,1]{label=\color{blue!75!black}\arabic*., font=\bfseries, leftmargin=2em, labelsep=1em}
\setlist{nosep} 

\title{Logique Mathématique}
\author{Mignard Mael}
\date{\today}

\begin{document}

\maketitle
\tableofcontents
\newpage

\section{Logique des propositions}
\subsection{Définitions}

C'est un énoncé déclaratif qu'on peut dire vrai ou faux, sinon ce n'est pas une proposition.
\begin{itemize}
    \item Exemples de propositions : "Il pleut", "2+2=4", "Le ciel est bleu".
    \item Exemple de prédicat (une "proposition" avec une variable) : "$x+1=4$".
    \item Exemples de non-propositions : "Quel temps fait-il ?", "Fais tes devoirs !", "Si j'étais riche...".
\end{itemize}

\subsection{Opérateurs logiques}

\begin{itemize}
    \item Négation : $\neg P$ (non P) est vrai si et seulement si P est faux.
    \item Conjonction : $P \wedge Q$ (P et Q) est vrai si et seulement si P et Q sont tous les deux vrais.
    \item Disjonction : $P \vee Q$ (P ou Q) est vrai si au moins l'un de P ou Q est vrai.
    \item Implication : $P \Rightarrow Q$ (si P alors Q) est faux uniquement lorsque P est vrai et Q est faux.
    \item Équivalence : $P \Leftrightarrow Q$ (P si et seulement si Q) est vrai lorsque P et Q ont la même valeur de vérité.
\end{itemize}

\subsection{Tableau de vérité des opérateurs logiques}

\begin{center}
\begin{tabular}{|c|c|c|c|c|c|}
\hline
P & Q & $\neg P$ & $P \wedge Q$ & $P \vee Q$ & $P \Rightarrow Q$ \\
\hline
V & V & F & V & V & V \\ 
V & F & F & F & V & F \\ 
F & V & V & F & V & V \\ 
F & F & V & F & F & V \\ 
\hline
\end{tabular}
\end{center}

\subsection{Exercices}

Construire les tableaux de vérité pour les expressions suivantes :

1) $P \wedge \neg Q$ 
\begin{center}
\begin{tabular}{|c|c|c|}
\hline
P & Q & $P \wedge \neg Q$ \\
\hline
V & V & F \\
V & F & V \\
F & V & F \\
F & F & F \\
\hline
\end{tabular}
\end{center}

\newpage
2) $\neg P \vee Q$ 
\begin{center}
\begin{tabular}{|c|c|c|}
\hline
P & Q & $\neg P \vee Q$ \\
\hline
V & V & F \\
V & F & V \\
F & V & F \\
F & F & F \\
\hline
\end{tabular}
\end{center}

3) $(P \wedge \neg Q) \vee (\neg P \wedge Q)$ 
\begin{center}
\begin{tabular}{|c|c|c|}
\hline
P & Q & $(P \wedge \neg Q) \vee (\neg P \wedge Q)$ \\
\hline
V & V & F \\
V & F & V \\
F & V & V \\
F & F & F \\
\hline
\end{tabular}
\end{center}

4) Conclure.
On a construit le ou exclusif (XOR) : $P \oplus Q = (P \wedge \neg Q) \vee (\neg P \wedge Q)$.

\section{Propriétés des formules}
\subsection{Tautologie}
Une formule est une tautologie si sa table de vérité ne contient que des "Vrai".
Exemple : $P \vee \neg P$ ("Être ou ne pas être"), toujours vrai.

\subsection{Contradiction}
Une formule est une contradiction si sa table de vérité ne contient que des "Faux".
Exemple : $P \wedge \neg P$, toujours faux.

\subsection{Formule Neutre}
Ni une tautologie, ni une contradiction (mélange de V et F).

\newpage
\section{Théorème de De Morgan}

\subsection{Négation d'une conjonction : $\neg (P \wedge Q)$}
La règle est : $\neg (P \wedge Q) \iff \neg P \vee \neg Q$. Vérification :
\begin{center}
\begin{tabular}{|c|c|c|c|}
\hline
P & Q & $P \wedge Q$ & $\neg (P \wedge Q)$ \\
\hline
F & F & F & V \\
F & V & F & V \\
V & F & F & V \\
V & V & V & F \\
\hline
\end{tabular}
\end{center}

\subsection{Négation d'une disjonction : $\neg (P \vee Q)$}
La règle est : $\neg (P \vee Q) \iff \neg P \wedge \neg Q$. Vérification :
\begin{center}
\begin{tabular}{|c|c|c|c|}
\hline
P & Q & $P \vee Q$ & $\neg (P \vee Q)$ \\
\hline
F & F & F & V \\
F & V & V & F \\ 
V & F & V & F \\
V & V & V & F \\
\hline
\end{tabular}
\end{center}

\newpage
\section{Quantificateurs}
\subsection{Les différents quantificateurs}

\subsubsection{Quantificateur universel : $\forall$}
$ \forall x, P(x) $ signifie que pour tout $x$, $P(x)$ est vrai.

\subsubsection{Quantificateur existentiel : $\exists$}
$ \exists x, P(x) $ signifie qu'il existe au moins un $x$ pour lequel $P(x)$ est vrai.

\subsubsection{Exercices d'application}
\begin{itemize}
    \item $ \forall x \in \mathbb{R}, x^2 \geq 0 $ (Vrai)
    \item $ \exists x \in \mathbb{N}, x^2 = 25 $ (Vrai, $x=5$)
    \item $ \forall x \in \mathbb{N}, x^2 = x $ (Faux, $2^2 \neq 2$)
    \item $ \exists x \in \mathbb{R}, x^2 = -1 $ (Faux)
    \item $ \exists x \in \mathbb{C}, x^2 = -1 $ (Vrai, $x=i$)
\end{itemize}

\subsection{Négation des quantificateurs}
\subsubsection{Règles de passage}
\begin{itemize}
    \item $\neg (\forall x, P(x)) \iff \exists x, \neg P(x)$
    \item $\neg (\exists x, P(x)) \iff \forall x, \neg P(x)$
\end{itemize}

\subsubsection{Exercices de négation}
Faisons la négation des propositions précédentes :
\begin{itemize}
    \item $ \exists x \in \mathbb{R}, x^2 < 0 $ (Faux)
    \item $ \forall x \in \mathbb{N}, x^2 \neq 25 $ (Faux)
    \item $ \exists x \in \mathbb{N}, x^2 \neq x $ (Vrai, $x=2$)
    \item $ \forall x \in \mathbb{R}, x^2 \neq -1 $ (Vrai)
    \item $\forall x \in \mathbb{C}, x^2 \neq -1$ (Faux)
\end{itemize}

\end{document}